\section{Methodology}\label{sec:evaluation-methodology}

The benchmark was run locally with a Release build of \texttt{motivoc}.
The binary was built outside the repository at \texttt{/tmp/motivo-thesis-bench-build} using CMake 4.3.2 and Apple
clang 21.0.0.
Measurements were taken on an Apple M3 Pro machine with 18 GB of RAM\@.
The timing tool was \texttt{hyperfine} 1.20.0 with shell execution disabled, 5 warmup runs, and 50 measured
runs~\cite{Hyperfine_2026}.
Each measured command compiled one \texttt{.motivo} source file and wrote the resulting MIDI file, so the
results include
process startup, parsing, semantic analysis, lowering, MIDI serialization, and file output.

The benchmark corpus was synthetic.
It was built to exercise source-to-event expansion rather than to produce musically refined pieces.
This distinction matters.
A loop that generates 20,000 notes is useful for measuring expansion and compilation cost, but it does not prove that
the generated score is artistically interesting.
The benchmark therefore supports the engineering claim that Motivo can generate many symbolic events from
compact source rules within the tested limit.

A defining characteristic of the Motivo compiler is its ``zero-runtime'' architecture, which shifts the
computational burden entirely to the compilation phase.
To evaluate the viability of this approach, it is essential to analyze the compiler's execution speed, particularly as
the complexity of the musical composition scales.

\subsection{Theoretical Complexity}\label{detail-subsec:theoretical-complexity}

The compilation pipeline is composed of several sequential passes, each with a distinct computational complexity:

\noindent \textbf{Frontend and Semantic Analysis.} The complexity of parsing and type-checking is proportional to the
number of lines of source code, $O(S)$, where $S$ is the size of the abstract syntax tree.
Because these passes do not unroll loops or instantiate patterns multiple times, they remain extremely fast
regardless of the composition's length.

\noindent \textbf{Lowering Engine.} The unrolling process has a time complexity of $O(E)$, where $E$ is the total number
of generated \texttt{NoteEvent}s. A simple \texttt{loop (1000)} statement will execute its body 1,000 times during
this phase, meaning the lowering time is proportional to the density of the final timeline, not the brevity of the
source code.

\noindent \textbf{MIDI Backend.} The backend must sort the absolute events by their tick values before computing
delta-times.
This stable sort introduces a time complexity of $O(E \log E)$.
Serializing the events into the binary file requires a final linear pass, $O(E)$.

Therefore, the overall time complexity of the compilation process is bounded by the sorting phase in the backend:
$O(E \log E)$.
